\section{Background}
\label{sec:background}

\paragraph{Output parameter.} A C function can return a value through a
parameter of pointer-to-pointer type, written \texttt{T **P}. The caller
passes the address of its own pointer variable; the function writes a
value into \texttt{*P} for the caller to read after the call returns. We
write \texttt{*P} for the object the caller's pointer refers to, not for
\texttt{P} itself.

\paragraph{Release and reset.} A \emph{release} of \texttt{*P} is a call
that frees the object \texttt{*P} points to, for example
\texttt{OPENSSL\_free(*P)}. After a release, \texttt{*P} still holds the
freed address unless something changes it. A \emph{reset} of \texttt{*P}
is an assignment that gives it a new value, most often \texttt{*P = NULL}.
A dangling \texttt{*P}, one that still holds a freed address after a
release and before any reset, is dangerous exactly when the caller can
read or free it again: a double free, or a use of freed memory if the
caller dereferences it.

\paragraph{Typestate.} \emph{Typestate analysis} tracks not just an
object's type but which operations are currently valid on it, as a small
state machine. \todo{Cite the originating typestate paper once its
author list and venue are verified directly (for example on IEEE Xplore);
automated search returned conflicting author lists for this title across
mirrors and we do not cite it unverified.} For the invariant in this paper,
\texttt{*P} moves through three states: \emph{Valid} before any release,
\emph{Released} after a release and before a reset, and back to
\emph{Valid} (we call it \emph{Reset}, to keep the release that caused it
visible in our output) once a reset occurs. A \emph{violation} is a
program path on which \texttt{*P} can reach function exit while still in
the \emph{Released} state. CVE-2022-4450 (Section~\ref{sec:rq1}) is
exactly this: \texttt{PEM\_read\_bio\_ex} frees \texttt{*header} and
\texttt{*data} on one failure path and returns without resetting either.

\paragraph{Release site.} We call the specific call expression that
performs a release, together with its source location and the parameter
\texttt{P} it releases, a \emph{release site}. Section~\ref{sec:eval}
reports counts of release sites, not of functions: one function can
contain several release sites, one per output parameter it releases.

\paragraph{Rule as data.} Existing patch-matching checkers, including the
Part~1 baseline in this project (\texttt{typestate\_check.py}), encode a
rule as a sequence of string or syntax checks written for one named
function. We instead write the rule once as data
(Section~\ref{sec:rule}), independent of any function name, and apply it
to every function in a source tree that has a parameter of the tracked
shape.
